\section{Which Flatness Matters Under Constrained Adaptation?}
\label{sec:which-flatness}

The classical flatness literature is largely inherited from full-parameter training, where perturbations and curvature are defined over the whole model.
In PECL, however, the loss remains task-conditioned while the admissible update directions are sharply restricted by the frozen-backbone mechanism.
This creates a basic conceptual question: once future adaptation is confined to $\mathcal{W}_{\Delta,t}$, which notion of flatness remains relevant to continual generalization and forgetting?

\subsection{Why This Question Arises in LoRA-Based Continual Learning}
\label{sec:flatness-question}

In full-parameter training, flatness is typically understood through perturbations in the ambient parameter space.
Sharpness-aware methods favor solutions whose loss remains stable under such local perturbations.
In LoRA-based continual learning, by contrast, only admissible low-rank update directions are available to fit new task data.
The loss is still defined on the full model, but the optimization mechanism can only traverse a highly constrained subset of directions.

This would already motivate a re-examination of flatness.
But the issue is sharpened by the empirical observations from the introduction.
Applying sharpness-aware optimization \emph{only} within the LoRA update subspace can already yield a solution whose surrounding landscape appears flatter even when probed in the full parameter space, while simultaneously improving continual performance.
Thus, flatness control imposed on admissible LoRA directions already has consequences beyond those directions---yet the mechanism governing this effect must be understood in terms of the constrained update geometry.

The resulting question is therefore not merely about implementation convenience.
It is a question of \emph{relevance under constrained adaptation}: when future updates are limited to admissible low-rank directions, should continual generalization be understood through ambient full-space flatness, or through task-conditioned local flatness along admissible update directions?

\subsection{Three Flatness Notions: Ambient, Coordinate, and Admissible}
\label{sec:three-flatness}

To answer this question, it is useful to distinguish three related but non-equivalent notions of flatness.

\paragraph{Ambient full-space flatness.}
This is the standard notion inherited from full-parameter training.
It measures local sensitivity of the loss with respect to perturbations in the entire parameter space $\mathbb{R}^d$, including both the frozen backbone and the adapter.
Classical sharpness-aware methods such as SAM~\citep{foretSharpnessAwareMinimizationEfficiently2021} and Flat-LoRA~\citep{li2024flatlora} are formulated at this level.

\paragraph{Trainable-coordinate flatness.}
This measures local sensitivity with respect to perturbations in the trainable parameter coordinates.
In LoRA, these are the factor coordinates $(A_t, B_t) \in \Theta_\Delta$.
This notion is implementation-dependent: it is defined in the coordinate space in which optimization is performed, not in weight-increment space.
LoRA-SAM~\citep{NEURIPS2024_4eb2c0ad} implicitly uses this notion by applying SAM directly to $(A_t, B_t)$.

\paragraph{Admissible-update flatness.}
This measures local sensitivity with respect to perturbations supported on the task-admissible update subspace $\mathcal{W}_{\Delta,t}$.
It is a geometric notion defined in weight-increment space, and it captures sensitivity along directions that the PECL mechanism can actually realize under the frozen-backbone constraint.
This is the flatness object studied by our theory.

\paragraph{Where the three coincide and where they diverge.}
These three notions coincide in full-parameter fine-tuning, where every direction is simultaneously ambient, trainable, and admissible.
In PECL they generally diverge.
The backbone is frozen, so the set of admissible future updates is strictly smaller than the ambient parameter space, while the LoRA factor coordinates provide only a local parameterization of the corresponding admissible increment geometry.
In this paper, \emph{trainable-coordinate flatness is treated as an implementation-level notion}, whereas \emph{admissible-update flatness is the geometric object of interest}: it is defined in weight-increment space and invariant to how the LoRA factors are parameterized locally.

\subsection{Why Inherited Full-Space Flatness Is Not Automatically Relevant}
\label{sec:why-subspace}

The key issue in PECL is not whether full-space flatness may correlate with good performance, but whether it is the quantity that future continual adaptation can actually exploit.

Consider a model after learning task $t$: $W_t = W_t^{\mathrm{froz}} + \Delta W_t$.
Under the frozen-backbone PECL constraint, the next task cannot modify $W_t$ arbitrarily in the ambient parameter space.
It can only move the model through admissible update directions supported on $\mathcal{W}_{\Delta,t+1}$.
This immediately creates a mismatch with flatness notions inherited from full-parameter training, which aggregate curvature over many directions that are unreachable under the PECL mechanism.

A model may therefore appear flat in the ambient full parameter space while still being sharp along the admissible update directions available to future tasks.
Such a solution is not genuinely favorable for PECL, because future adaptation must still traverse sharp barriers in the only directions it can realize.
Conversely, a model may remain sensitive in many frozen directions while being flat along the admissible update geometry---and would be well-suited to continual adaptation despite imperfect ambient flatness.

This identifies the central conceptual shift.
In frozen-backbone PECL, the relevant notion of flatness is task-conditioned, solution-local flatness supported on the admissible update geometry---not arbitrary global flatness in the ambient parameter space.
Full-space flatness may remain a useful empirical correlate (and later we will see that shaping admissible-update flatness can alter the surrounding full-model landscape as a byproduct), but full-space flatness is not the directly theoretically governed object once adaptation is constrained to admissible low-rank directions.

\textbf{What this paper does and does not prove.}
We prove that admissible-update flatness is the \emph{bound-relevant} notion in frozen-backbone PECL, under the conditions that (i) the backbone is strictly frozen, (ii) posteriors are supported on $\mathcal{W}_{\Delta,t}$, and (iii) the Jacobian bridge identifies $\mathcal{W}_{\Delta,t}$ with the local first-order reachable directions.
We do not prove a general equivalence between ambient and admissible-update flatness, nor do we claim that ambient flatness is always harmful.
Violations of condition~(i)---such as partially unfreezing the backbone---would weaken the reduction.
We test this boundary condition in Section~\ref{sec:boundary}.

\subsection{Sharpness Measures Under Admissible Geometry}
\label{section:sahrpness}

We now introduce the sharpness measures used throughout the paper.
Motivated by PAC-Bayesian analysis, SAM~\citep{foretSharpnessAwareMinimizationEfficiently2021} minimizes the \textbf{adversarial zeroth-order sharpness}:
\[
\mathrm{AS}^{(0)}_S(W,\rho) = \max_{\|\epsilon\| \leq \rho} \bigl[\hat{L}_S(W+\epsilon) - \hat{L}_S(W)\bigr].
\]
A complementary formulation is the \textbf{random zeroth-order sharpness}~\citep{pmlr-v151-bisla22a,jiang2019fantasticgeneralizationmeasures}:
\[
\mathrm{RS}_S^{(0)}(W, \Sigma) = \mathbb{E}_{\varepsilon \sim \mathcal{N}(0,\Sigma)} \bigl[\hat{L}_S(W+\varepsilon) - \hat{L}_S(W)\bigr].
\]
The \textbf{adversarial first-order sharpness} of GAM~\citep{Zhang_GAM} targets the worst-case gradient norm:
\[
\mathrm{AS}^{(1)}_S(W,\rho) = \max_{\|\epsilon\| \leq \rho} \|\nabla_W \hat{L}_S(W+\epsilon)\|.
\]
For any $W$ and radius $\rho$, $\mathrm{AS}^{(1)}_S$ upper-bounds $\mathrm{AS}^{(0)}_S$, making first-order sharpness a stronger flatness criterion~\citep{Zhang_GAM}.

In full-parameter training, these measures are typically interpreted in the ambient parameter space.
In PECL, however, all sharpness quantities considered in this paper are interpreted relative to the admissible update geometry: we replace $W$ with $W_t^{\mathrm{froz}} + \widehat{\Delta W}_t$ and restrict perturbations $\epsilon$ to directions supported on $\mathcal{W}_{\Delta,t}$.
This gives the PECL-specific quantities $\mathrm{AS}^{(0),\Delta}$, $\mathrm{RS}^{(0),\Delta}$, and $\mathrm{AS}^{(1),\Delta}$ used in the experiments.
The next section will ask not merely how sharpness is measured, but which sharpness-dependent terms actually enter the continual generalization bound under frozen-backbone PECL.
